% Test of listings too tall to keep together (didactic): a listing of at most
% keeplines lines that doesn't fit on a page must break, as a long listing
% does, with a warning, and nothing may run off the paper.
%
% Build: see the Makefile, which checks this document's PDF and log:
%   - no "Overfull \vbox";
%   - every word of every listing is in the PDF (pdftotext leaves out what
%     is off the paper);
%   - the warnings below, each with its input line.
%
% The options are explicit, so that the cases don't depend on the defaults.
% test-keep-tall100.tex is the same with keeplines=100, for the cases of a
% long listing that the option asks to keep.
%
% This document, keeplines=20:
%   T1  20 lines of minted[breaklines], each wrapped over many rows: pages
%       of them.  Warns "too tall".  1600 times "word".
%   T3  6 lines of more than a third of a page each, so that keepends=3 lines
%       don't fit on a page: breaks between them anyway.  Warns "too tall".
%       2041 times "third", with the lead-in.
%   T4  A long listing (24 lines) of such lines: warns "can't keep".
%       8161 times "long", with the lead-in.
%   T5  Something between the lines that can't be taken apart (a \mark, by
%       a redefined \FV@ListProcessLine): warns "something between", all lines
%       stay.
%       8 times "marked".
% test-keep-tall100.tex, keeplines=100:
%   T2  90 lines: warns "too tall".  90 times "unit".
%   T6  A listing that fits the text block, but not with the heading and
%       lead-in before it: breaks, warns "too tall", and no overfull page.
\documentclass[a4paper]{memoir}
\usepackage[british]{babel}
\usepackage{amsmath,amssymb,amsthm}
\ifdefined\TALLOPTS
  \PassOptionsToPackage{keeplines=100}{didactic}
\else
  \PassOptionsToPackage{keeplines=20}{didactic}
\fi
\usepackage{didactic}
\usepackage[noamsthm,notheorems]{beamerarticle}
\makeatletter
\newwrite\@talltest@out
\newcount\@talltest@i
\newcount\@talltest@j
% \talltest{name}{lines}{words per line}{word}: write \jobname-name.txt.
\newcommand\talltest[4]{%
  \immediate\openout\@talltest@out=\jobname-#1.txt
  \def\@talltest@line{}%
  \@talltest@j\z@
  \@whilenum\@talltest@j<#3\do{%
    \advance\@talltest@j\@ne
    \edef\@talltest@line{\@talltest@line #4 }}%
  \@talltest@i\z@
  \@whilenum\@talltest@i<#2\do{%
    \advance\@talltest@i\@ne
    \immediate\write\@talltest@out{\@talltest@line}}%
  \immediate\closeout\@talltest@out}
\newcount\talltestfit
\makeatother
\begin{document}
\ifdefined\TALLOPTS
\section{T2}
Lead-in to a listing of 90 lines, which the option asks to keep:
\talltest{u}{90}{1}{unit}%
\VerbatimInput{\jobname-u.txt}
Text after T2.

\clearpage
\section{T6}
Lead-in to a listing that fits the text block, but not with this:
% one line short of the text block, in lines of the text
\talltestfit=\numexpr\textheight/\baselineskip-1\relax
\talltest{f}{\talltestfit}{1}{fitting}%
\VerbatimInput{\jobname-f.txt}
Text after T6.
\else
\section{T1}
Lead-in to a listing of 20 lines that wrap:
\talltest{a}{20}{80}{word}%
\inputminted[breaklines]{text}{\jobname-a.txt}
Text after T1.

\clearpage
\section{T3}
Lead-in to lines of a third of a page:
\talltest{c}{6}{340}{third}%
\VerbatimInput[breaklines,fontsize=\footnotesize]{\jobname-c.txt}
Text after T3.

\clearpage
\section{T4}
Lead-in to a long listing of such lines:
\talltest{d}{24}{340}{long}%
\VerbatimInput[breaklines,fontsize=\footnotesize]{\jobname-d.txt}
Text after T4.

\clearpage
\section{T5}
Lead-in to a listing with something between its lines:
\talltest{e}{8}{1}{marked}%
\begingroup
\makeatletter
\let\talltestline\FV@ListProcessLine
\def\FV@ListProcessLine#1{\talltestline{#1}\mark{}}%
\makeatother
\VerbatimInput{\jobname-e.txt}
\endgroup
Text after T5.
\fi
\end{document}
